\documentclass[a4paper,10pt]{article}
\usepackage[utf8]{inputenc}
\usepackage{graphicx}
\usepackage{subfigure}
\usepackage{fancyhdr}
\usepackage{listings}
\usepackage{ gensymb }
\usepackage{hyperref}
\usepackage{multicol}


%opening

\pagestyle{fancy} %eigener Seitenstil
\fancyhf{} %alle Kopf- und Fußzeilenfelder bereinigen
\fancyhead[L]{SEP Abgabegespräch SS2015/16} %Kopfzeile links
\fancyhead[C]{} %zentrierte Kopfzeile
\fancyhead[R]{\includegraphics[width=1cm]{../../template/tugraz_bw}} %Kopfzeile rechts
\renewcommand{\headrulewidth}{0.4pt} %obere Trennlinie
\fancyfoot[C]{} %Seitennummer
\renewcommand{\footrulewidth}{0.4pt} %untere Trennlinie


\lstset{language=C++, basicstyle=\ttfamily\tiny}  

\begin{document}

\section*{Snake}

Entwickeln Sie eine vereinfachte Version des Spieles Snake\\ (Bsp.: \url{http://playsnake.org/}).
Wie im Basisbeispiel soll ein zweidimensionales Spielfeld angelegt werden. Das Spielfeld ist dabei
10x10 Felder groß. Die Spalten sollen für die Ausgabe mit einem Leerzeichen getrennt werden.\\
Jedes Feld wird Anfangs mit einem Leerzeichen initialisiert mit Ausnahme der Felder
am äußersten Rand. Die Randfelder sollen eine Wand mit dem Zeichen \texttt{\#} (ASCII 35) darstellen.\\
Der Schlangenkörper wird mit dem Zeichen \texttt{O} (ASCII 79) gefüllt, wobei der Schlangenkopf mit dem Zeichen 
\texttt{X} (ASCII 88) dargestellt wird. Die Bewegungsbefehle sind die gleichen wie beim Basisbeispiel:\\
\texttt{u} (ASCII 117) für nach oben, \texttt{d} (ASCII 100) für nach unten, \texttt{l} (ASCII 108) für nach links und 
\texttt{r} (ASCII 114) für nach rechts. Mit dem Befehl \texttt{quit} soll das Spiel beendet werden.\\
Die Größe und Startposition der Schlange soll wie in der ersten Ausgabe des Spielfeldes auf der Konsole sein.
Es gibt keine Gewinnbedingung. Bei Auftreten von spielinternen Fehlern soll Nichts geschehen.


\begin{multicols}{2}

\subsection*{Beispiel}
\begin{verbatim}
Input auf der Konsole: 
sep> d
sep> r
sep> quit
\end{verbatim}

\subsubsection*{Ausgabe}

\begin{verbatim}
Konsolen Ausgabe:
# # # # # # # # # # 
#                 # 
#                 # 
#   O O X         # 
#                 # 
#                 # 
#                 # 
#                 # 
#                 # 
# # # # # # # # # # 
# # # # # # # # # # 
#                 # 
#                 # 
#     O O         # 
#       X         # 
#                 # 
#                 # 
#                 # 
#                 # 
# # # # # # # # # # 
# # # # # # # # # # 
#                 # 
#                 # 
#       O         # 
#       O X       # 
#                 # 
#                 # 
#                 # 
#                 # 
# # # # # # # # # # 
\end{verbatim}

\end{multicols}


\subsection*{Aufruf}
\texttt{./snake}

\subsection*{Bewertung}
\begin{itemize}
 \item Dokumentation und Programmierstil: 1 Punkt
 \item Einhalten der OOP Konzepte: 1 Punkt
 \item Struktur \& Korrektheit: 1 Punkt
\end{itemize}

\newpage


\end{document}
